\documentclass[12pt,a4paper]{article}

% ── Packages ─────────────────────────────────────────────────────────────────
\usepackage[a4paper, margin=2.5cm]{geometry}
\usepackage{fontenc}
\usepackage{inputenc}
\usepackage{lmodern}
\usepackage{microtype}
\usepackage{amsmath, amssymb}
\usepackage{graphicx}
\usepackage{xcolor}
\usepackage{listings}
\usepackage{booktabs}
\usepackage{array}
\usepackage{longtable}
\usepackage{enumitem}
\usepackage{hyperref}
\usepackage{fancyhdr}
\usepackage{titlesec}
\usepackage{tcolorbox}
\usepackage{mdframed}
\usepackage{caption}
\usepackage{subcaption}
\usepackage{parskip}

% ── Colors ───────────────────────────────────────────────────────────────────
\definecolor{codebg}{RGB}{245,245,250}
\definecolor{codecomment}{RGB}{100,100,100}
\definecolor{codestring}{RGB}{30,120,30}
\definecolor{codekeyword}{RGB}{0,80,180}
\definecolor{codenumber}{RGB}{150,50,150}
\definecolor{noteblue}{RGB}{220,235,255}
\definecolor{noteborder}{RGB}{70,130,200}
\definecolor{warnbg}{RGB}{255,245,220}
\definecolor{warnborder}{RGB}{200,140,0}
\definecolor{titlecolor}{RGB}{30,60,120}
\definecolor{sectioncolor}{RGB}{0,80,160}

% ── Listings (code style) ─────────────────────────────────────────────────────
\lstdefinestyle{python}{
  language=Python,
  backgroundcolor=\color{codebg},
  basicstyle=\ttfamily\small,
  keywordstyle=\color{codekeyword}\bfseries,
  stringstyle=\color{codestring},
  commentstyle=\color{codecomment}\itshape,
  numberstyle=\tiny\color{codenumber},
  numbers=left,
  numbersep=8pt,
  stepnumber=1,
  frame=single,
  framerule=0.4pt,
  rulecolor=\color{gray!40},
  breaklines=true,
  showstringspaces=false,
  tabsize=4,
  captionpos=b,
  xleftmargin=1.5em,
}

\lstdefinestyle{json}{
  language={},
  backgroundcolor=\color{codebg},
  basicstyle=\ttfamily\small,
  stringstyle=\color{codestring},
  keywordstyle=\color{codekeyword},
  frame=single,
  framerule=0.4pt,
  rulecolor=\color{gray!40},
  breaklines=true,
  showstringspaces=false,
  tabsize=2,
  xleftmargin=1.5em,
}

\lstdefinestyle{bash}{
  language=bash,
  backgroundcolor=\color{codebg},
  basicstyle=\ttfamily\small,
  keywordstyle=\color{codekeyword},
  frame=single,
  framerule=0.4pt,
  rulecolor=\color{gray!40},
  breaklines=true,
  showstringspaces=false,
  xleftmargin=1.5em,
}

% ── Note/Warning boxes ────────────────────────────────────────────────────────
\newmdenv[
  backgroundcolor=noteblue,
  linecolor=noteborder,
  linewidth=2pt,
  leftline=true, rightline=false, topline=false, bottomline=false,
  innerleftmargin=10pt,
  innerrightmargin=10pt,
  innertopmargin=8pt,
  innerbottommargin=8pt,
  skipabove=\baselineskip,
  skipbelow=\baselineskip,
]{notebox}

\newmdenv[
  backgroundcolor=warnbg,
  linecolor=warnborder,
  linewidth=2pt,
  leftline=true, rightline=false, topline=false, bottomline=false,
  innerleftmargin=10pt,
  innerrightmargin=10pt,
  innertopmargin=8pt,
  innerbottommargin=8pt,
  skipabove=\baselineskip,
  skipbelow=\baselineskip,
]{warnbox}

% ── Section formatting ────────────────────────────────────────────────────────
\titleformat{\section}{\large\bfseries\color{sectioncolor}}{{\thesection}}{1em}{}[\titlerule]
\titleformat{\subsection}{\normalsize\bfseries\color{sectioncolor!80}}{\thesubsection}{1em}{}
\titleformat{\subsubsection}{\normalsize\bfseries}{\thesubsubsection}{1em}{}

% ── Header / footer ───────────────────────────────────────────────────────────
\pagestyle{fancy}
\fancyhf{}
\fancyhead[L]{\small\textcolor{gray}{CUDA-METRO User Guide}}
\fancyhead[R]{\small\textcolor{gray}{\leftmark}}
\fancyfoot[C]{\small\thepage}
\renewcommand{\headrulewidth}{0.4pt}

% ── Hyperlinks ────────────────────────────────────────────────────────────────
\hypersetup{
  colorlinks=true,
  linkcolor=sectioncolor,
  urlcolor=sectioncolor!80,
  citecolor=sectioncolor,
}

% ── Macros ────────────────────────────────────────────────────────────────────
\newcommand{\code}[1]{\texttt{#1}}
\newcommand{\file}[1]{\texttt{\textit{#1}}}
\newcommand{\param}[1]{\texttt{#1}}

% ─────────────────────────────────────────────────────────────────────────────

\begin{document}

% ── Title page ───────────────────────────────────────────────────────────────
\begin{titlepage}
  \centering
  \vspace*{2cm}
  {\Huge\bfseries\color{titlecolor} CUDA-METRO}\\[0.5em]
  {\large\color{gray} GPU-Accelerated Metropolis Monte Carlo\\for 2D Magnetic Materials}\\[2.5cm]

  {\Large\bfseries User Guide}\\[0.5em]
  {\normalsize A practical introduction for Python programmers}\\[3cm]

  \begin{mdframed}[
    backgroundcolor=noteblue,
    linecolor=noteborder,
    linewidth=1pt,
    innerleftmargin=20pt,
    innerrightmargin=20pt,
    innertopmargin=15pt,
    innerbottommargin=15pt,
  ]
  \centering
  \textbf{About this guide}\\[0.5em]
  \begin{flushleft}
  This guide assumes you are comfortable writing Python scripts, using
  \code{numpy} arrays, and running programs from the terminal.
  It does \emph{not} assume prior knowledge of GPU programming, CUDA, or
  condensed matter physics. Physics concepts are explained as needed.
  \end{flushleft}
  \end{mdframed}

  \vfill
  {\small Based on: Hait \& Mahapatra (2025), \textit{Journal of Open Source Software}, 10(109), 7589}\\[0.5em]
  {\small\color{gray} CUDA-METRO --- \url{https://github.com/arkavo/CUDA-METRO}}
\end{titlepage}

\tableofcontents
\newpage

% ─────────────────────────────────────────────────────────────────────────────
\section{What Is CUDA-METRO?}
% ─────────────────────────────────────────────────────────────────────────────

CUDA-METRO is a Python library for simulating the magnetic behaviour of
thin 2D materials on an NVIDIA GPU. It answers questions like:

\begin{itemize}
  \item At what temperature does a magnet lose its ordering (Curie temperature)?
  \item What spin texture forms when a material is cooled slowly?
  \item How does the magnetisation respond when you sweep an external field?
\end{itemize}

The simulation runs a \textbf{Metropolis Monte Carlo} algorithm: spins are
proposed random moves and accepted or rejected based on energy. Running
millions of such steps lets the system relax to a physically realistic state.
Doing this on a GPU instead of a CPU makes it 100--1000$\times$ faster.

\subsection{High-level architecture}

\begin{center}
\begin{tabular}{lll}
\toprule
\textbf{File} & \textbf{Language} & \textbf{Role} \\
\midrule
\file{construct.py} & Python & User-facing API; orchestrates the simulation \\
\file{montecarlo.py} & Python + CUDA C & GPU kernels compiled at import time \\
\file{hysteresis.py} & Python & $M_z$--$H$ hysteresis experiment \\
\file{multilayer.py} & Python & Multilayer equilibration with field zones \\
\file{visualize.py} & Python & Post-run spin visualisation \\
\file{inputs/TC\_*.csv} & CSV & Material parameter database ($>$60 materials) \\
\bottomrule
\end{tabular}
\end{center}

The two classes you will use are:
\begin{itemize}
  \item \code{MonteCarlo} --- creates and runs a simulation.
  \item \code{Analyze} --- loads saved results and generates plots.
\end{itemize}

% ─────────────────────────────────────────────────────────────────────────────
\section{Physics Background (Short Version)}
% ─────────────────────────────────────────────────────────────────────────────

You do not need to be a physicist to use this software, but a little
background helps you understand what the parameters mean.

\subsection{Spins on a lattice}

A magnetic material is modelled as a 2D grid of atoms. Each atom carries a
\textbf{spin} --- a tiny magnetic moment pointing in some direction in 3D
space. We represent it as a 3-component vector:
$\mathbf{S}_i = (S_x, S_y, S_z)$.

The magnitude $|\mathbf{S}_i|$ equals the \textbf{spin quantum number} $S$
(e.g.\ $S = 3/2$ for Cr$^{3+}$ ions in CrPS$_4$). At each MC step we
propose a new random direction for one spin and decide whether to accept it.

\subsection{The Hamiltonian (energy function)}

The energy of spin $i$ interacting with its neighbours is:

\begin{equation}
  H_i = -\sum_j J_j\, \mathbf{S}_i \cdot \mathbf{S}_j
        - \sum_{j,\alpha} K_j^\alpha S_i^\alpha S_j^\alpha
        - A \, S_{iz}^2
        - \sum_j \mathbf{D}_{ij} \cdot (\mathbf{S}_i \times \mathbf{S}_j)
        - \mu\, B \cdot S_{iz}
  \label{eq:hamiltonian}
\end{equation}

\begin{center}
\begin{tabular}{clp{8cm}}
\toprule
\textbf{Term} & \textbf{Symbol} & \textbf{Physical meaning} \\
\midrule
Exchange & $J_j$ & Tendency for neighbours to align ($J>0$: ferromagnetic,
                    all spins point the same way; $J<0$: antiferromagnetic,
                    alternating) \\
Anisotropic exchange & $K_j^\alpha$ & Direction-dependent version of exchange \\
Single-ion anisotropy & $A$ & Preference for spins to point along a particular
                              axis (easy-axis if $A<0$, easy-plane if $A>0$) \\
DMI & $\mathbf{D}_{ij}$ & Dzyaloshinskii-Moriya interaction; causes spins to
                           twist at an angle, stabilising skyrmions \\
External field & $B$ & Applied magnetic field; pulls $S_z$ toward alignment \\
\bottomrule
\end{tabular}
\end{center}

\subsection{Metropolis Monte Carlo}

Monte Carlo simulation generates a sequence of spin configurations whose
statistics match thermodynamic equilibrium at temperature $T$. The
\textbf{Metropolis algorithm} achieves this with one simple rule per
proposed move:

\begin{enumerate}
  \item Pick a spin $i$ at random.
  \item Generate a trial spin direction $\mathbf{S}'_i$ uniformly on the unit sphere.
  \item Compute the energy change $\Delta E = H(\text{trial}) - H(\text{current})$.
  \item Accept the trial if $\Delta E < 0$ (lower energy is always better), \emph{or} with probability $e^{-\Delta E / k_B T}$ even if $\Delta E > 0$ (thermal fluctuations can push the system uphill).
\end{enumerate}

Repeating this millions of times lets the system find low-energy spin
configurations. At high temperature, thermal energy is large so many uphill
moves are accepted and the spins are disordered. At low temperature, almost
only downhill moves are accepted and order builds up.

\begin{notebox}
\textbf{Why GPU?} --- Each proposed move only touches a small neighbourhood of
spins, so many proposals can be evaluated \emph{simultaneously} on different
parts of the lattice. An NVIDIA GPU can run thousands of such threads in
parallel, making the simulation orders of magnitude faster than CPU code.
\end{notebox}

% ─────────────────────────────────────────────────────────────────────────────
\section{Installation and Dependencies}
% ─────────────────────────────────────────────────────────────────────────────

\subsection{Requirements}

\begin{itemize}
  \item NVIDIA GPU with CUDA support (any compute capability $\geq 5.0$).
  \item CUDA Toolkit installed (version 11 or newer recommended).
  \item Python 3.9+.
  \item The following Python packages:
\end{itemize}

\begin{center}
\begin{tabular}{ll}
\toprule
\textbf{Package} & \textbf{Purpose} \\
\midrule
\code{pycuda} & Python bindings for CUDA; compiles and runs GPU kernels \\
\code{numpy} & Array operations on the CPU side \\
\code{matplotlib} & Plotting \\
\code{seaborn} & Heatmap visualisation \\
\code{tqdm} & Progress bars \\
\bottomrule
\end{tabular}
\end{center}

\subsection{Installing packages}

\begin{lstlisting}[style=bash]
pip install pycuda numpy matplotlib seaborn tqdm
\end{lstlisting}

\begin{warnbox}
\textbf{pyCUDA requires a working CUDA installation.} If \code{pip install pycuda}
fails, check that \code{nvcc --version} works in your terminal. You may need to
set \code{CUDA\_ROOT} or install the CUDA toolkit from NVIDIA's website first.
\end{warnbox}

\subsection{Running scripts}

All scripts are run from inside the \file{src/cudametro/} directory:

\begin{lstlisting}[style=bash]
cd src/cudametro
python sample.py
\end{lstlisting}

% ─────────────────────────────────────────────────────────────────────────────
\section{Configuration File}
% ─────────────────────────────────────────────────────────────────────────────

Every simulation is controlled by a \textbf{JSON configuration file}. Here is
a fully annotated example:

\begin{lstlisting}[style=json, caption={Annotated configuration file}]
{
    "Single_Mat_Flag" : 1,      // 1 = use a single material (default)
    "Animation_Flags" : 0,      // 0 = do not save animation frames
    "DMI_Flag"        : 1,      // 1 = include Dzyaloshinskii-Moriya term
    "TC_Flag"         : 0,      // 1 = sweep temperature to find Tc
    "Static_T_Flag"   : 1,      // 1 = fixed temperature; 0 = temperature sweep
    "FM_Flag"         : 0,      // 1 = start all spins pointing up (+z)
                                //     0 = start with random spin directions
    "Input_flag"      : 0,      // 1 = load a saved spin state as starting point
    "Input_File"      : "Starting_States/state.npy",

    "Temps"           : ["2.0"],     // temperature(s) in Kelvin
    "Material"        : "CrPS4",     // material name (must match inputs/TC_*.csv)

    "SIZE"    : 200,            // lattice side length; grid is SIZE x SIZE
    "Box"     : 16,             // internal tiling parameter
    "Blocks"  : 2048,           // parallel MC proposals per inner loop step
    "Threads" : 2,              // threads per block (always 2; do not change)

    "B"       : "0.0",          // initial external magnetic field (normalised)

    "stability_runs"   : 300000,  // inner loop length: sweeps per wrap pass
    "calculation_runs" : 50,      // sweeps used for measurement averaging
    "stability_wrap"   : 30,      // outer loop: how many wrap passes to run
    "calculation_wrap" : 35,

    "Cmpl_Flag" : 1,            // 1 = complex (full anisotropic) Hamiltonian
    "Prefix"    : "CrPS4"       // prefix for the output directory name
}
\end{lstlisting}

\subsection{The most important parameters}

\subsubsection{SIZE}
The lattice is a \param{SIZE}~$\times$~\param{SIZE} grid. A 200$\times$200
grid has 40\,000 sites. Larger lattices capture long-range structures (e.g.\
large skyrmions) but use more GPU memory and take longer.

\subsubsection{Blocks}
\param{Blocks} is the number of independent spin proposals evaluated
\emph{simultaneously} on the GPU per inner-loop iteration. Larger values
speed up each iteration but consume more GPU memory (the pre-allocated
random number arrays scale as \code{Blocks} $\times$ \code{stability\_runs}).

A good rule of thumb is $\texttt{Blocks} \approx 0.1 \times
\texttt{SIZE}^2$ (10\% of all sites).

\subsubsection{stability\_runs and stability\_wrap}
The outer loop runs \param{stability\_wrap} times. Each outer iteration runs
\param{stability\_runs} inner Metropolis sweeps and saves one grid snapshot.
Total sweeps = \param{stability\_wrap}~$\times$~\param{stability\_runs}.

\subsubsection{FM\_Flag}
\begin{itemize}
  \item \param{FM\_Flag = 1}: All spins start pointing along $+z$. Good for
        studying ferromagnets at low field.
  \item \param{FM\_Flag = 0}: Spins start in random directions. Required when
        applying field zones (§\ref{sec:fieldzones}) to avoid the system being
        so stable that no moves are accepted.
\end{itemize}

% ─────────────────────────────────────────────────────────────────────────────
\section{Material Parameters}
% ─────────────────────────────────────────────────────────────────────────────

Material properties are stored in CSV files in the \file{inputs/} folder. The
filename convention is \file{TC\_MaterialName.csv}. The config field
\code{"Material": "CrPS4"} instructs the code to load
\file{inputs/TC\_CrPS4.csv}.

\subsection{CSV column order}

\begin{center}
\begin{tabular}{cclr}
\toprule
\textbf{Array index} & \textbf{Symbol} & \textbf{Description} & \textbf{Unit} \\
\midrule
0  & $S$         & Spin quantum number               & --- \\
1  & $J_1$       & 1st-neighbour exchange            & meV \\
2  & $J_2$       & 2nd-neighbour exchange            & meV \\
3  & $J_3$       & 3rd-neighbour exchange            & meV \\
4  & $J_4$       & 4th-neighbour exchange            & meV \\
5--7  & $K_{1x,y,z}$ & Anisotropic exchange, shell 1 & meV \\
8--10 & $K_{2x,y,z}$ & Anisotropic exchange, shell 2 & meV \\
11--13 & $K_{3x,y,z}$ & Anisotropic exchange, shell 3 & meV \\
14--16 & $K_{4x,y,z}$ & Anisotropic exchange, shell 4 & meV \\
17--19 & $A_{x,y,z}$ & Single-ion anisotropy         & meV \\
20--23 & NBS       & Neighbour-shell descriptor        & int \\
24 & $T_c$        & Reference Curie temperature        & K \\
25 & $D$          & DMI magnitude                     & meV \\
\bottomrule
\end{tabular}
\end{center}

\subsection{Available materials}

The \file{inputs/} directory includes over 60 parameterised materials, including
CrI$_3$, CrBr$_3$, CrCl$_3$, CrPS$_4$, CrPS$_4$-ML, FePS$_3$, MnSe$_2$,
VSe$_2$, and many more. Use \code{"Material": "CrI3"} in your config to load
\file{inputs/TC\_CrI3.csv}.

\subsection{The neighbour structure descriptor (NBS)}

The NBS field (indices 20--23) encodes how many neighbours exist at each shell
distance. The code uses this to select the correct CUDA kernel:

\begin{center}
\begin{tabular}{lll}
\toprule
\textbf{NBS} & \textbf{Crystal} & \textbf{Typical materials} \\
\midrule
3-6-3-6  & Honeycomb & CrI$_3$, CrPS$_4$, CrBr$_3$ \\
6-6-6-12 & Triangular & CrTe$_2$, MnSe$_2$ \\
4-4-4-8  & Square    & CrO$_2$, MnO$_2$ \\
2-2-4-2  & Rectangular & --- \\
2-4-2-4  & Rectangular variant & --- \\
\bottomrule
\end{tabular}
\end{center}

% ─────────────────────────────────────────────────────────────────────────────
\section{Running a Simulation (Single Layer)}
\label{sec:singlelayer}
% ─────────────────────────────────────────────────────────────────────────────

\subsection{The minimal workflow}

\begin{lstlisting}[style=python, caption={Minimal single-layer simulation --- sample.py}]
import construct as cst
import numpy as np
from tqdm import tqdm

# 1. Create the simulation object from a config file
sim = cst.MonteCarlo(config="sample_config.json")

# 2. Initialise the spin grid
sim.mc_init()
sim.display_material()   # print material parameters to terminal

# 3. Run the equilibration loop
for i in tqdm(range(sim.S_Wrap), desc="Stability Runs"):
    sim.generate_random_numbers(sim.stability_runs)
    grid = sim.run_mc_dmi_66612(sim.T[0])   # returns current grid
    np.save(f"{sim.save_directory}/grid_{i:04d}", grid)

print("Done. Results saved in:", sim.save_directory)
\end{lstlisting}

\subsection{Step-by-step explanation}

\paragraph{Step 1: \code{MonteCarlo(config=...)}}
Reads the JSON config, loads the material CSV, allocates all GPU memory, and
creates a timestamped output directory (e.g.\ \file{Output\_CrPS4\_2025\_01\_15\_14\_30\_00/}).
A \file{metadata.json} summary is written there automatically.

\paragraph{Step 2: \code{mc\_init()}}
Allocates the spin grid on the CPU and GPU. If \param{FM\_Flag=1}, all spins
are initialised to $(S_x, S_y, S_z) = (0, 0, S)$. If \param{FM\_Flag=0}, each
spin is placed at a random point on the unit sphere and scaled by $S$.

\paragraph{Step 3: The outer loop}
Each iteration of the loop (controlled by \param{stability\_wrap}) does:
\begin{enumerate}
  \item \code{generate\_random\_numbers}: fill four GPU arrays with fresh uniform
        random floats (for site selection, trial spin direction, and
        Metropolis acceptance test).
  \item \code{run\_mc\_*}: run \param{stability\_runs} Metropolis sweeps on the
        GPU. Returns the updated spin grid as a flat \code{numpy} array.
  \item \code{np.save}: write the current grid to disk for later visualisation.
\end{enumerate}

\subsection{Choosing the right \code{run\_mc\_*} method}

Pick the method that matches your material's crystal structure:

\begin{center}
\begin{tabular}{ll}
\toprule
\textbf{Method} & \textbf{Crystal / Use} \\
\midrule
\code{run\_mc\_dmi\_66612(T)} & Triangular lattice (6-6-6-12) with DMI \\
\code{run\_mc\_dmi\_4448(T)}  & Square lattice (4-4-4-8) with DMI \\
\code{run\_mc\_dmi\_3636(T)}  & Honeycomb lattice (3-6-3-6) with DMI \\
\code{run\_mc\_3636\_multilayer(T, layers)} & Honeycomb multilayer \\
\bottomrule
\end{tabular}
\end{center}

\subsection{Understanding the output directory}

After the run, the output folder contains:
\begin{itemize}
  \item \file{metadata.json} --- all simulation parameters.
  \item \file{grid\_0000.npy}, \file{grid\_0001.npy}, \dots ---
        one spin grid snapshot per \param{stability\_wrap} iteration.
        Each is a flat \code{float32} array of length
        $\texttt{SIZE}^2 \times 3$.
\end{itemize}

% ─────────────────────────────────────────────────────────────────────────────
\section{The Spin Grid: Memory Layout}
\label{sec:gridlayout}
% ─────────────────────────────────────────────────────────────────────────────

Understanding the grid layout is essential when you want to analyse or plot
results yourself.

\subsection{Single-layer layout}

\begin{lstlisting}[style=python]
# grid is a 1D float32 array of length SIZE*SIZE*3
grid = sim.grid

# Index of spin component k at site (row, col):
# k=0 -> Sx,  k=1 -> Sy,  k=2 -> Sz
site = row * SIZE + col
value = grid[site * 3 + k]

# Or reshape to a 3D array for convenient access:
g = grid.reshape(SIZE, SIZE, 3)
Sz = g[:, :, 2]   # 2D array of z-components
\end{lstlisting}

\subsection{Multilayer layout}

For $L$ layers the grid is extended along a new outermost axis:

\begin{lstlisting}[style=python]
# grid is a 1D float32 array of length L*SIZE*SIZE*3
# Layer index is outermost:
pti = layer * SIZE*SIZE  +  row*SIZE  +  col
value = grid[pti * 3 + k]

# Or reshape to a 4D array:
g = grid.reshape(L, SIZE, SIZE, 3)
Sz_layer2 = g[2, :, :, 2]   # z-component of layer 2
\end{lstlisting}

Flat memory is laid out as:

\begin{center}
\begin{tabular}{|c|c|c|c|}
\hline
Layer 0 ($N^2 \times 3$ floats) & Layer 1 & Layer 2 & \ldots Layer $L{-}1$ \\
\hline
\end{tabular}
\end{center}

% ─────────────────────────────────────────────────────────────────────────────
\section{Multilayer Simulations}
\label{sec:multilayer}
% ─────────────────────────────────────────────────────────────────────────────

A multilayer system stacks $L$ 2D honeycomb planes on top of each other.
Spins within each layer interact via the usual in-plane exchange and DMI terms.
Adjacent layers are coupled through an \textbf{interlayer exchange}
$J_c$ (stored in the CSV file).

\subsection{Physical picture}

\begin{center}
\begin{tabular}{c}
\ttfamily\small
\begin{tabular}{|c|}
\hline
$\cdot\cdot\cdot\cdot$ Layer $L-1$ (top) \\
\hline
$\cdot\cdot\cdot\cdot$ Layer $L-2$ \\
\hline
$\quad\vdots\quad$ \\
\hline
$\cdot\cdot\cdot\cdot$ Layer 1 \\
\hline
$\cdot\cdot\cdot\cdot$ Layer 0 (bottom) \\
\hline
\end{tabular}
\end{tabular}
\end{center}

Each site in layer $\ell > 0$ feels an additional energy contribution from the
spin directly below it (same row and column, layer $\ell - 1$) and above it
(layer $\ell + 1$):
\[
H_{\text{inter}} = J_c \, \mathbf{S}_{i,\ell} \cdot \mathbf{S}_{i,\ell-1}
                 + J_c \, \mathbf{S}_{i,\ell} \cdot \mathbf{S}_{i,\ell+1}
\]

For CrPS$_4$-ML, $J_c = -0.16\,\text{meV}$ (antiferromagnetic), meaning
adjacent layers prefer to point in opposite directions. The competition
between the ferromagnetic in-plane coupling and the antiferromagnetic
interlayer coupling produces rich magnetic phase diagrams.

\subsection{Running a multilayer simulation}

\begin{lstlisting}[style=python, caption={multilayer.py workflow}]
import construct as cst
import numpy as np
from tqdm import tqdm

LAYERS = 8

sim = cst.MonteCarlo(config="multilayer_config.json")
sim.mc_init_multilayer(layers=LAYERS)    # note: different init call
sim.display_material()

# (Optional) apply a spatially varying field -- see Section 8
sim.set_field_zones(
    boundaries=[10, 50, 250, 290],
    fields=[0.0, 0.0, 0.6, 0.0, 0.0]
)

for i in tqdm(range(sim.S_Wrap), desc="Stability Runs"):
    sim.generate_random_numbers_multilayer(None)      # multilayer version
    grid = sim.run_mc_3636_multilayer(sim.T[0], layers=LAYERS)
    np.save(f"{sim.save_directory}/grid_{i:04d}", grid)
\end{lstlisting}

The key differences from the single-layer workflow are:
\begin{enumerate}
  \item Use \code{mc\_init\_multilayer(layers=L)} instead of \code{mc\_init()}.
  \item Use \code{generate\_random\_numbers\_multilayer(None)} (the argument is
        unused but kept for API consistency).
  \item Use \code{run\_mc\_3636\_multilayer(T, layers=L)}.
\end{enumerate}

\subsection{Spin initialisation modes}

\begin{center}
\begin{tabular}{lp{10cm}}
\toprule
\textbf{FM\_Flag} & \textbf{Starting state} \\
\midrule
1 & \textbf{Ferromagnetic}: every spin in every layer is set to
    $(S_x, S_y, S_z) = (0, 0, S)$ --- all pointing up. The system is
    strongly stable at low temperature; very few moves are accepted
    initially. \\[0.3em]
0 & \textbf{Random}: each spin is placed at a uniformly random point
    on the unit sphere. The system starts disordered and must equilibrate.
    Required when applying field zones so that the spins have room to
    respond to the spatial field gradient. \\
\bottomrule
\end{tabular}
\end{center}

\subsection{Measuring magnetisation}

\begin{lstlisting}[style=python]
# Returns (M_total, M_per_layer)
M_tot, M_layers = sim.get_magnetization(layers=LAYERS)

# M_tot: scalar, mean Sz / spin S, over all layers and sites
# M_layers: numpy array of shape (LAYERS,), per-layer average
print(f"Total Mz/Msat = {M_tot:.4f}")
for i, ml in enumerate(M_layers):
    print(f"  Layer {i}: Mz/Msat = {ml:.4f}")
\end{lstlisting}

Only the $z$-component of the spin ($S_z$) is used, normalised by the spin
magnitude $S$. This gives $M_z / M_\text{sat} \in [-1, +1]$.

% ─────────────────────────────────────────────────────────────────────────────
\section{Spatially Varying Magnetic Field Zones}
\label{sec:fieldzones}
% ─────────────────────────────────────────────────────────────────────────────

By default, a single uniform external field is applied to the entire lattice.
CUDA-METRO also supports \textbf{field zones}: different columns of the
lattice can have different field values.

\subsection{How it works}

The GPU stores the field as an array of \param{SIZE} floats, one value per
column. The Hamiltonian reads \code{b[col]} for each site, where
\code{col = pti \% SIZE}. This means field zones are strips running vertically
(parallel to the row axis).

\subsection{API}

\begin{lstlisting}[style=python]
# Uniform field (backward-compatible):
sim.set_field(0.6)

# Spatially varying field:
# boundaries: list of column indices dividing the lattice into zones
# fields: one value per zone (must have len(boundaries) + 1 entries)
sim.set_field_zones(
    boundaries=[10, 50, 250, 290],   # 4 boundaries -> 5 zones
    fields=[0.0, 0.0, 0.6, 0.0, 0.0]
)
\end{lstlisting}

\subsection{Zone layout example}

For a 300$\times$300 lattice with \code{boundaries=[10, 50, 250, 290]}:

\begin{center}
\begin{tabular}{|c|c|c|c|c|}
\hline
\textbf{Zone 0} & \textbf{Zone 1} & \textbf{Zone 2} & \textbf{Zone 3} & \textbf{Zone 4} \\
cols 0--9 & cols 10--49 & cols 50--249 & cols 250--289 & cols 290--299 \\
$B = 0.0$ & $B = 0.0$ & $B = 0.6$ & $B = 0.0$ & $B = 0.0$ \\
(buffer) & (buffer) & (active) & (buffer) & (buffer) \\
\hline
\end{tabular}
\end{center}

The outer buffer zones have zero field and act as a magnetically neutral
boundary. The central zone carries the experiment field. This design prevents
edge effects from corrupting the region of interest.

\begin{warnbox}
\textbf{Always use \code{FM\_Flag = 0} (random start) with field zones.}
Starting from an all-up FM state with a positive $z$-field makes the initial
configuration so energetically stable that the acceptance rate collapses to
near zero and the simulation appears frozen.
\end{warnbox}

\subsection{Number of zones is flexible}

The number of zones is not locked to 5. Provide any list of boundaries and a
matching \code{fields} list:

\begin{lstlisting}[style=python]
# 3 zones (two buffers, one central)
sim.set_field_zones(
    boundaries=[50, 150],
    fields=[0.0, 1.0, 0.0]
)

# 7 zones
sim.set_field_zones(
    boundaries=[20, 50, 100, 200, 250, 280],
    fields=[0.0, 0.2, 0.4, 0.6, 0.4, 0.2, 0.0]
)
\end{lstlisting}

% ─────────────────────────────────────────────────────────────────────────────
\section{The Hysteresis Experiment}
\label{sec:hysteresis}
% ─────────────────────────────────────────────────────────────────────────────

\file{hysteresis.py} measures the $M_z$--$H$ hysteresis loop: how the
magnetisation responds as the external field is swept up, then down, then back
up again. This is the same measurement made in a physical magnetometer.

\subsection{Field sweep protocol}

The field is swept in three branches:

\[
0 \;\to\; +H_\text{max} \;\to\; -H_\text{max} \;\to\; +H_\text{max}
\]

\begin{lstlisting}[style=python]
h1 = np.linspace(0,       H_MAX,  N_STEPS,    endpoint=False)  # ramp up
h2 = np.linspace(H_MAX,  -H_MAX, 2*N_STEPS,   endpoint=False)  # down sweep
h3 = np.linspace(-H_MAX,  H_MAX, 2*N_STEPS,   endpoint=True)   # up sweep
H_sweep = np.concatenate([h1, h2, h3])
\end{lstlisting}

At each field value:
\begin{enumerate}
  \item \code{set\_field(h\_val)} updates the GPU field array.
  \item \code{EQ\_WRAP} passes of \code{generate + run\_mc} equilibrate the
        spins at this field.
  \item \code{get\_magnetization()} records $M_z / M_\text{sat}$ and
        $M_z / M_\text{sat}$ per layer.
\end{enumerate}

Crucially, the spin state is \textbf{not reset} between field steps. The
configuration at field $H_n$ seeds the simulation at $H_{n+1}$. This is what
produces hysteresis: the system remembers which way it was magnetised.

\subsection{Equilibration (EQ\_WRAP)}

\param{EQ\_WRAP} is the number of \code{generate + run\_mc} passes executed at
each field point. More passes allow the spins more time to respond to a field
change. The total number of Metropolis sweeps at each field point is:

\[
\text{sweeps per field point} = \texttt{EQ\_WRAP} \times \texttt{stability\_runs}
\]

At low temperatures (below $\sim 5\,\text{K}$ for CrPS$_4$-ML) thermal
fluctuations are rare and the system needs more sweeps to overcome local energy
barriers. Increasing \param{EQ\_WRAP} does not use extra GPU memory; it only
adds compute time.

\subsection{Key parameters}

\begin{center}
\begin{tabular}{lll}
\toprule
\textbf{Parameter} & \textbf{Default} & \textbf{Meaning} \\
\midrule
\code{LAYERS}       & 12     & Number of magnetic layers \\
\code{TEMPERATURES} & [0.1, 0.5, 1, 2, 5, 10, 15, 20]\,K & Temperatures to simulate \\
\code{H\_MAX}       & 1.5    & Maximum normalised field \\
\code{N\_STEPS}     & 4      & Field points per branch \\
\code{EQ\_WRAP}     & 10     & Equilibration passes per field point \\
\bottomrule
\end{tabular}
\end{center}

\subsection{Outputs}

All outputs are saved in a timestamped directory
\file{Hysteresis\_YYYY\_MM\_DD\_HH\_MM\_SS/}:

\begin{center}
\begin{tabular}{lp{9cm}}
\toprule
\textbf{File} & \textbf{Content} \\
\midrule
\file{hysteresis\_data.npy} & Raw data dictionary; keys are temperatures;
                               each entry contains arrays \code{H}, \code{M\_total},
                               \code{M\_layers}, and \code{snapshots} \\
\file{MH\_all\_temps.png}   & $M_z/M_\text{sat}$ and emu/g vs $H$ for all
                               temperatures; line style encodes branch \\
\file{MH\_per\_layer\_T\{X\}K.png} & Per-layer $M_z$ vs $H$ at each temperature \\
\file{spinstate\_T\{X\}K\_H0.png} & Heatmaps of $S_x$, $S_y$, $S_z$ components
                                     at $H \approx 0$ for each temperature \\
\bottomrule
\end{tabular}
\end{center}

\subsection{Magnetisation units}

Magnetisation is reported in two units side by side:

\begin{itemize}
  \item \textbf{Normalised}: $M_z / M_\text{sat} \in [-1, +1]$, where $+1$
        means all spins fully aligned along $+z$.
  \item \textbf{emu/g}: the physical unit used in magnetometry. For CrPS$_4$:
        \[
          M_\text{sat}^{\text{emu/g}} = \frac{g \cdot \mu_B \cdot S \cdot N_A}{M_\text{molar}}
          = \frac{2.0 \times 9.274\times10^{-21} \times 1.5 \times 6.022\times10^{23}}{211.21}
          \approx 79.3\,\text{emu/g}
        \]
\end{itemize}

\subsection{Loading saved data}

\begin{lstlisting}[style=python]
import numpy as np

data = np.load("Hysteresis_2025_01_15_14_30_00/hysteresis_data.npy",
               allow_pickle=True).item()

# data is a dict keyed by temperature (float)
T = 2.0
H = data[T]["H"]           # field sweep array
M = data[T]["M_total"]     # total Mz/Msat at each field point
Ml = data[T]["M_layers"]   # shape (N_H, LAYERS): per-layer Mz
snaps = data[T]["snapshots"]  # {H_value: grid_array}
\end{lstlisting}

% ─────────────────────────────────────────────────────────────────────────────
\section{Visualisation and Analysis}
\label{sec:viz}
% ─────────────────────────────────────────────────────────────────────────────

\subsection{The Analyze class}

\begin{lstlisting}[style=python]
viewer = cst.Analyze(
    directory="Output_CrPS4_2025_01_15_14_30_00",
    reverse=False,     # False = plot frames in forward order
)
\end{lstlisting}

\subsubsection{Single-layer methods}

\begin{lstlisting}[style=python]
viewer.spin_view()       # heatmaps of Sx, Sy, Sz for each saved grid
viewer.quiver_view()     # in-plane (Sx, Sy) arrows coloured by Sz
\end{lstlisting}

\subsubsection{Multilayer methods}

\begin{lstlisting}[style=python]
viewer.spin_view_multilayer(layers=8)    # 3 x L heatmap grid
viewer.quiver_view_multilayer(layers=8)  # quiver per layer, shared Sz colour
\end{lstlisting}

\subsection{Manual analysis example}

\begin{lstlisting}[style=python]
import numpy as np
import matplotlib.pyplot as plt

SIZE = 300
LAYERS = 8

# Load the last saved grid
grid = np.load("Output_CrPS4/grid_0029.npy")
g = grid.reshape(LAYERS, SIZE, SIZE, 3)

fig, axes = plt.subplots(1, LAYERS, figsize=(4*LAYERS, 4))
for i, ax in enumerate(axes):
    ax.imshow(g[i, :, :, 2], cmap="coolwarm", vmin=-1.5, vmax=1.5)
    ax.set_title(f"Layer {i}")
    ax.axis("off")
plt.suptitle("Sz per layer")
plt.tight_layout()
plt.savefig("sz_layers.png")
\end{lstlisting}

% ─────────────────────────────────────────────────────────────────────────────
\section{Understanding CUDA Concepts (as a Python User)}
\label{sec:cuda}
% ─────────────────────────────────────────────────────────────────────────────

You do not need to write any CUDA code yourself. But understanding what happens
behind the scenes helps you tune performance and avoid common mistakes.

\subsection{GPU vs CPU memory}

A GPU has its own memory (VRAM), separate from the computer's regular RAM. Data
must be explicitly copied between them. In pyCUDA:

\begin{lstlisting}[style=python]
# CPU -> GPU (host to device)
drv.memcpy_htod(GPU_array, cpu_numpy_array)

# GPU -> CPU (device to host)
drv.memcpy_dtoh(cpu_numpy_array, GPU_array)
\end{lstlisting}

CUDA-METRO handles all of this for you inside \code{MonteCarlo}. When you call
\code{run\_mc\_*()}, the function launches the GPU kernel, lets it run, then
copies the result back to \code{sim.grid} on the CPU.

\subsection{CUDA kernels}

A CUDA kernel is a function that runs on the GPU, executed simultaneously by
many threads. Each thread handles one spin proposal. The two threads in each
block in CUDA-METRO work as a pair: one computes the energy of the current
spin, the other computes the energy of the trial spin. They synchronise,
compare, and write the result.

\subsection{Pre-allocated random numbers}

Generating random numbers on the GPU is expensive. CUDA-METRO pre-allocates
eight large GPU arrays at \code{mc\_init\_multilayer} time:

\[
\text{size of each array} = \texttt{Blocks} \times \texttt{stability\_runs}
\]

Each outer loop iteration calls \code{rg.fill\_uniform()} to refill these in
place (no new allocation). This is why changing \param{Blocks} or
\param{stability\_runs} after initialisation is not supported --- the array
sizes are fixed at init time.

\subsection{VRAM budget}

\begin{lstlisting}[style=python]
import pycuda.driver as drv
free, total = drv.mem_get_info()
used_mb = (total - free) / 1024**2
print(f"VRAM used: {used_mb:.1f} MB")
\end{lstlisting}

The dominant allocation is the random number arrays. For
\param{Blocks = 4096}, \param{stability\_runs = 60000}:
\[
  8 \;\text{arrays} \times 4096 \times 60000 \times 4\,\text{bytes}
  \approx 7.5\,\text{GB}
\]
On a GPU with 32\,GB VRAM, you have room to push this further.

% ─────────────────────────────────────────────────────────────────────────────
\section{Common Mistakes and How to Avoid Them}
\label{sec:mistakes}
% ─────────────────────────────────────────────────────────────────────────────

\begin{enumerate}

\item \textbf{Simulation appears frozen / no spin changes}

  \textit{Cause}: \code{FM\_Flag = 1} combined with a field in the $+z$
  direction. The system is fully polarised and every trial spin
  (random direction) has higher energy --- almost no moves are accepted.

  \textit{Fix}: Set \code{FM\_Flag = 0} in the config to start from a random
  state, especially when using field zones or studying antiferromagnets.

\item \textbf{Wrong number of layers initialised}

  If you call \code{mc\_init\_multilayer(layers=8)} but your config was designed
  for 4 layers, the spin grid will have 8 layers but the saved metadata will say
  4. Always keep \code{LAYERS} consistent between the Python script and the config
  file.

\item \textbf{VRAM out-of-memory error}

  \textit{Cause}: \param{Blocks} $\times$ \param{stability\_runs} is too large.

  \textit{Fix}: Reduce one of them. A useful diagnostic:
  \begin{lstlisting}[style=python]
import pycuda.driver as drv
free, total = drv.mem_get_info()
print(f"Free VRAM: {free/1e9:.1f} GB")
  \end{lstlisting}

\item \textbf{Magnetisation does not saturate at low temperature}

  \textit{Cause}: Too few equilibration passes. At very low $T$, the system
  needs more sweeps to tunnel out of local energy minima.

  \textit{Fix}: Increase \param{EQ\_WRAP} (hysteresis) or \param{stability\_wrap}
  (multilayer) at the cost of longer run time.

\item \textbf{Wrong material loaded}

  The config field \code{"Material"} must exactly match the CSV filename stem
  (case-sensitive on Linux). For \file{inputs/TC\_CrI3.csv}, use
  \code{"Material": "CrI3"}.

\item \textbf{ImportError when running outside \file{src/cudametro/}}

  \code{construct.py} uses relative imports. Always run scripts from
  \file{src/cudametro/} or add the directory to \code{sys.path}:
  \begin{lstlisting}[style=python]
import sys
sys.path.insert(0, "/path/to/src/cudametro")
import construct as cst
  \end{lstlisting}

\end{enumerate}

% ─────────────────────────────────────────────────────────────────────────────
\section{Performance Tuning}
\label{sec:perf}
% ─────────────────────────────────────────────────────────────────────────────

\subsection{Parallelism ratio}

The fraction of sites updated per inner step is:
\[
  f = \frac{\texttt{Blocks}}{\texttt{SIZE}^2}
\]

Higher $f$ speeds up convergence per wall-clock second but reduces accuracy (the
strict Metropolis detailed-balance condition is violated more). A ratio of
$f \approx 0.1$ (10\% of sites per step) is the recommended default, giving
accuracy comparable to single-spin-update algorithms.

\subsection{Benchmark reference points}

\begin{center}
\begin{tabular}{llr}
\toprule
\textbf{Lattice} & \textbf{GPU} & \textbf{Time per full run} \\
\midrule
200$\times$200 & NVIDIA V100 & $\sim$30\,s \\
500$\times$500 & NVIDIA A100 & $\sim$300\,s \\
750$\times$750 & NVIDIA A100 & $\sim$9\,h \\
\bottomrule
\end{tabular}
\end{center}

\subsection{Tuning for a specific GPU}

If you have a GPU with large VRAM (e.g.\ 32\,GB), you can increase
\param{stability\_runs} and \param{Blocks} to run more sweeps per outer
iteration without saving to disk as often. This reduces I/O overhead for very
long runs.

% ─────────────────────────────────────────────────────────────────────────────
\section{Quick-Reference Tables}
\label{sec:ref}
% ─────────────────────────────────────────────────────────────────────────────

\subsection{MonteCarlo API}

\begin{center}
\begin{longtable}{lp{10cm}}
\toprule
\textbf{Method} & \textbf{Description} \\
\midrule
\endhead
\code{MonteCarlo(config)} & Constructor. Reads config, loads material, allocates GPU memory. \\
\code{mc\_init()} & Initialise single-layer spin grid. \\
\code{mc\_init\_multilayer(layers=L)} & Initialise multilayer spin grid. \\
\code{display\_material()} & Print material parameters to terminal. \\
\code{generate\_random\_numbers(n)} & Single-layer: generate RNG arrays. \\
\code{generate\_random\_numbers\_multilayer(n)} & Multilayer: refill pre-allocated RNG arrays. \\
\code{run\_mc\_dmi\_66612(T)} & Run MC for triangular (6-6-6-12) crystal. \\
\code{run\_mc\_dmi\_4448(T)} & Run MC for square (4-4-4-8) crystal. \\
\code{run\_mc\_dmi\_3636(T)} & Run MC for honeycomb (3-6-3-6) crystal. \\
\code{run\_mc\_3636\_multilayer(T, layers)} & Run MC for multilayer honeycomb system. \\
\code{set\_field(B)} & Broadcast a uniform field to all columns. \\
\code{set\_field\_zones(boundaries, fields)} & Set per-column spatially varying field. \\
\code{get\_magnetization(layers)} & Return $(M_z/M_\text{sat},\; M_z\text{ per layer})$. \\
\code{grid\_reset()} & Re-initialise the spin grid from scratch. \\
\code{cleanup()} & Free GPU buffers. \\
\bottomrule
\end{longtable}
\end{center}

\subsection{Config key summary}

\begin{center}
\begin{longtable}{llp{7cm}}
\toprule
\textbf{Key} & \textbf{Type} & \textbf{Effect} \\
\midrule
\endhead
\code{SIZE}              & int    & Lattice side length ($N \times N$) \\
\code{Blocks}            & int    & Parallel MC proposals per step \\
\code{stability\_runs}   & int    & Sweeps per outer loop iteration \\
\code{stability\_wrap}   & int    & Number of outer loop iterations \\
\code{FM\_Flag}          & 0/1    & 0=random start; 1=all-up FM start \\
\code{DMI\_Flag}         & 0/1    & Include DMI term in Hamiltonian \\
\code{TC\_Flag}          & 0/1    & 1=temperature sweep (Tc mode) \\
\code{Static\_T\_Flag}   & 0/1    & 1=fixed T; 0=linspace to 2$T_c$ \\
\code{Temps}             & list   & List of temperatures as strings \\
\code{Material}          & string & Stem of CSV file in inputs/ \\
\code{B}                 & string & Initial uniform field \\
\code{Prefix}            & string & Output directory prefix \\
\bottomrule
\end{longtable}
\end{center}

% ─────────────────────────────────────────────────────────────────────────────
\section{Worked Example: CrPS$_4$-ML Hysteresis}
\label{sec:example}
% ─────────────────────────────────────────────────────────────────────────────

This example walks through \file{hysteresis.py} line by line.

\subsection{Setup}

\begin{lstlisting}[style=python]
# Physical constants for emu/g conversion
MU_B_CGS = 9.2741e-21   # erg/G (CGS Bohr magneton)
N_A      = 6.02214e23
G_LANDE  = 2.0
SPIN     = 1.5           # Cr3+ has S = 3/2
M_MOLAR  = 51.996 + 30.974 + 4*32.06   # CrPS4 = 211.21 g/mol
EMU_G_SAT = G_LANDE * MU_B_CGS * SPIN * N_A / M_MOLAR  # ~79.3 emu/g
\end{lstlisting}

\subsection{Field sweep construction}

\begin{lstlisting}[style=python]
H_MAX   = 1.5    # maximum field (in units of spin magnitude)
N_STEPS = 4      # field points per branch

h1 = np.linspace(0,      H_MAX,  N_STEPS,    endpoint=False)
h2 = np.linspace(H_MAX, -H_MAX, 2*N_STEPS,   endpoint=False)
h3 = np.linspace(-H_MAX, H_MAX, 2*N_STEPS,   endpoint=True)
H_sweep = np.concatenate([h1, h2, h3])
# Total: N_STEPS + 2*N_STEPS + 2*N_STEPS = 5*N_STEPS = 20 field points
\end{lstlisting}

\subsection{Temperature loop}

\begin{lstlisting}[style=python]
TEMPERATURES = [0.1, 0.5, 1.0, 2.0, 5.0, 10.0, 15.0, 20.0]  # K
EQ_WRAP      = 10   # equilibration passes per field point

for T in TEMPERATURES:
    sim = cst.MonteCarlo(config="hysteresis_config.json")
    sim.mc_init_multilayer(layers=LAYERS)   # fresh random state each T

    M_total_list, M_layers_list = [], []
    snapshots = {}

    for idx, h_val in enumerate(H_sweep):
        sim.set_field(h_val)

        # Equilibrate at this field
        for _ in range(EQ_WRAP):
            sim.generate_random_numbers_multilayer(None)
            sim.run_mc_3636_multilayer(T, layers=LAYERS)

        # Measure
        M_tot, M_lay = sim.get_magnetization(LAYERS)
        M_total_list.append(M_tot)
        M_layers_list.append(M_lay.copy())

        # Save a snapshot at H~0
        if idx in SNAP_INDICES:
            snapshots[h_val] = sim.grid.copy()
\end{lstlisting}

\begin{notebox}
\textbf{Why a fresh instance per temperature?}
Each temperature sweep creates a new \code{MonteCarlo} object. This resets the
spin state to a fresh random configuration so that different temperatures are
independent experiments. The GPU memory from the previous instance will be
freed when the Python object goes out of scope.
\end{notebox}

% ─────────────────────────────────────────────────────────────────────────────
\section{Next Steps and Further Reading}
% ─────────────────────────────────────────────────────────────────────────────

\subsection{Running your own material}

\begin{enumerate}
  \item Find or compute the exchange parameters $J_1, J_2, \ldots$, DMI $D$,
        and anisotropy $A$ for your material (typically from DFT calculations
        or literature).
  \item Create a CSV file in \file{inputs/} following the format in
        Section~4.
  \item Set \code{"Material": "YourMaterial"} in the config.
  \item Adjust \code{"SIZE"} and \code{"Blocks"} for your GPU.
  \item Run \file{sample.py} or write your own driver script.
\end{enumerate}

\subsection{Finding the Curie temperature}

Set \code{"TC\_Flag": 1} and \code{"Static\_T\_Flag": 0} in the config. The
code will sweep temperature from near 0 to $2 T_c$ and return $\langle M \rangle$
and the susceptibility $\chi$ at each temperature. The Curie temperature is
identified as the peak of $\chi$.

\subsection{Cite the software}

If you use CUDA-METRO in a publication, please cite:

\begin{quote}
Hait \& Mahapatra (2025). \textit{CUDA-METRO: A GPU-Accelerated Metropolis
Monte Carlo Simulator for 2D Magnetic Materials.} Journal of Open Source
Software, 10(109), 7589.
\end{quote}

\subsection{Reporting issues}

Bugs and feature requests can be filed at the project's GitHub repository.

\end{document}
